% --- Archon physics formalization source begin ---
% archon:physics
% archon:covers PhyXMiniProblems/problem_phyx_mini_0833.lean
% archon:source-report reports/phyx_mini/problem_phyx_mini_0833.source.json

\chapter{Physics problem phyx\_mini\_0833}
\label{ch:PhyXMiniProblems_problem_phyx_mini_0833}

\paragraph{Problem source.}
\#\# Question

What is the magnitude of the strength of the electric field at the position indicated by the dot?

\#\# Answer choices

- A: A: 1.3 \textbackslash{}times 10\textasciicircum{}\{5\}N/C

- B: B: 5.3 \textbackslash{}times 10\textasciicircum{}\{5\}N/C

- C: C: 1.0 \textbackslash{}times 10\textasciicircum{}\{5\}N/C

- D: D: 1.34 \textbackslash{}times 10\textasciicircum{}\{5\}N/C

\#\# Figure caption (auxiliary; use the image as primary evidence)

The image depicts a diagram with three charged particles arranged on a rectangle. 

1. **Charges**:
   - The top left charge is a negative charge labeled "-5.0 nC" and shown in teal.
   - The top right charge is a positive charge labeled "10 nC" and shown in red.
   - The bottom right charge is also a positive charge labeled "10 nC" and shown in red.

2. **Distances**:
   - The distance between the negative charge and the top right positive charge is labeled "4.0 cm."
   - The vertical distance from the negative charge to a point below is "2.0 cm."

3. **Geometrical Layout**:
   - The arrangement forms a rectangle with dashed lines representing connections between the charges and the geometric points.
   - The bottom left point, marked by a black dot, serves as a reference point for the rectangle's dimensions.

The arrangement visually illustrates the relative positions and charges of particles likely for an electrostatic analysis.

\#\# Dataset metadata

Electrostatics; Physical Model Grounding Reasoning, Multi-Formula Reasoning

\paragraph{Recorded answer/context.}
D: D: 1.34 \textbackslash{}times 10\textasciicircum{}\{5\}N/C

\paragraph{Figure/image path.}
/root/proposal\_for\_physic/science-mango/phyx\_mini\_run/phyx\_data/test\_image/833.png

\paragraph{Formalization target.}
create a compiling Lean file with sorry bodies at `PhyXMiniProblems/problem\_phyx\_mini\_0833.lean`. The Lean declarations must preserve the physical quantities, dimensions or dimensional roles, figure labels, governing-law hypotheses, and final relation expressed by this problem.
Use Mathlib/Physlib names found through LeanExplore where available. If a physics API is missing, introduce faithful local abstractions rather than scalar placeholder aliases.

\begin{theorem}[Physics formalization target]
\label{thm:physics:phyx_mini_0833:target}
\lean{PhyXMiniProblems.ProblemPhyXMini0833.electricFieldMagnitude_answer_D}
\uses{def:physics:phyx-mini-0833:phyxminiproblems-problemphyxmini0833-rectangularpointchargesetup, def:physics:phyx-mini-0833:phyxminiproblems-problemphyxmini0833-matchesproblemandprimaryfigure, def:physics:phyx-mini-0833:phyxminiproblems-problemphyxmini0833-usessireferencedata, def:physics:phyx-mini-0833:phyxminiproblems-problemphyxmini0833-hasphysicalparameters, def:physics:phyx-mini-0833:phyxminiproblems-problemphyxmini0833-satisfiespointchargecoulomblawandsuperposition, def:physics:phyx-mini-0833:phyxminiproblems-problemphyxmini0833-isuniqueclosestdisplayedchoice}
The assigned autoformalize agent should translate this physics problem into Lean declarations in the covered file, with theorem and lemma proof bodies written as `by sorry`.
\end{theorem}
\begin{proof}
This is an autoformalization task, not a proof task. Produce statements that can later be proved by the physics prover without weakening the physical model.
\end{proof}

% --- Archon declaration topology begin ---
\section{Declaration topology}
\begin{definition}[Length Quantity]
  \label{def:physics:phyx-mini-0833:phyxminiproblems-problemphyxmini0833-lengthquantity}
  \lean{PhyXMiniProblems.ProblemPhyXMini0833.LengthQuantity}
  A nonnegative physical length
\end{definition}
\begin{proof}
  This is the declared model definition of Length Quantity.
\end{proof}
\begin{definition}[Signed Charge Quantity]
  \label{def:physics:phyx-mini-0833:phyxminiproblems-problemphyxmini0833-signedchargequantity}
  \lean{PhyXMiniProblems.ProblemPhyXMini0833.SignedChargeQuantity}
  A signed physical electric charge
\end{definition}
\begin{proof}
  This is the declared model definition of Signed Charge Quantity.
\end{proof}
\begin{definition}[electric Field Strength Dimension]
  \label{def:physics:phyx-mini-0833:phyxminiproblems-problemphyxmini0833-electricfieldstrengthdimension}
  \lean{PhyXMiniProblems.ProblemPhyXMini0833.electricFieldStrengthDimension}
  The physical dimension M L T⁻² C⁻¹ of electric-field strength
\end{definition}
\begin{proof}
  This is the declared model definition of electric Field Strength Dimension.
\end{proof}
\begin{definition}[Planar Position Quantity]
  \label{def:physics:phyx-mini-0833:phyxminiproblems-problemphyxmini0833-planarpositionquantity}
  \lean{PhyXMiniProblems.ProblemPhyXMini0833.PlanarPositionQuantity}
  A unit-independent signed planar position
\end{definition}
\begin{proof}
  This is the declared model definition of Planar Position Quantity.
\end{proof}
\begin{definition}[Planar Electric Field Quantity]
  \label{def:physics:phyx-mini-0833:phyxminiproblems-problemphyxmini0833-planarelectricfieldquantity}
  \lean{PhyXMiniProblems.ProblemPhyXMini0833.PlanarElectricFieldQuantity}
  \uses{def:physics:phyx-mini-0833:phyxminiproblems-problemphyxmini0833-electricfieldstrengthdimension}
  A unit-independent planar electric-field vector
\end{definition}
\begin{proof}
  This is the declared model definition of Planar Electric Field Quantity.
\end{proof}
\begin{definition}[x Axis]
  \label{def:physics:phyx-mini-0833:phyxminiproblems-problemphyxmini0833-xaxis}
  \lean{PhyXMiniProblems.ProblemPhyXMini0833.xAxis}
  Coordinate 0, horizontal and positive to the right in the image
\end{definition}
\begin{proof}
  This is the declared model definition of x Axis.
\end{proof}
\begin{definition}[y Axis]
  \label{def:physics:phyx-mini-0833:phyxminiproblems-problemphyxmini0833-yaxis}
  \lean{PhyXMiniProblems.ProblemPhyXMini0833.yAxis}
  Coordinate 1, vertical and positive upward in the image
\end{definition}
\begin{proof}
  This is the declared model definition of y Axis.
\end{proof}
\begin{definition}[length In Meters]
  \label{def:physics:phyx-mini-0833:phyxminiproblems-problemphyxmini0833-lengthinmeters}
  \lean{PhyXMiniProblems.ProblemPhyXMini0833.lengthInMeters}
  \uses{def:physics:phyx-mini-0833:phyxminiproblems-problemphyxmini0833-lengthquantity}
  Read a physical length in coherent-SI metres
\end{definition}
\begin{proof}
  This is the declared model definition of length In Meters.
\end{proof}
\begin{definition}[charge In Coulombs]
  \label{def:physics:phyx-mini-0833:phyxminiproblems-problemphyxmini0833-chargeincoulombs}
  \lean{PhyXMiniProblems.ProblemPhyXMini0833.chargeInCoulombs}
  \uses{def:physics:phyx-mini-0833:phyxminiproblems-problemphyxmini0833-signedchargequantity}
  Read a signed physical charge in coherent-SI coulombs
\end{definition}
\begin{proof}
  This is the declared model definition of charge In Coulombs.
\end{proof}
\begin{definition}[position Vector In Meters]
  \label{def:physics:phyx-mini-0833:phyxminiproblems-problemphyxmini0833-positionvectorinmeters}
  \lean{PhyXMiniProblems.ProblemPhyXMini0833.positionVectorInMeters}
  \uses{def:physics:phyx-mini-0833:phyxminiproblems-problemphyxmini0833-planarpositionquantity}
  Read a planar position as Cartesian coordinates in metres
\end{definition}
\begin{proof}
  This is the declared model definition of position Vector In Meters.
\end{proof}
\begin{definition}[electric Field Vector In Newtons Per Coulomb]
  \label{def:physics:phyx-mini-0833:phyxminiproblems-problemphyxmini0833-electricfieldvectorinnewtonspercoulomb}
  \lean{PhyXMiniProblems.ProblemPhyXMini0833.electricFieldVectorInNewtonsPerCoulomb}
  \uses{def:physics:phyx-mini-0833:phyxminiproblems-problemphyxmini0833-planarelectricfieldquantity}
  Read a planar electric-field vector in newtons per coulomb
\end{definition}
\begin{proof}
  This is the declared model definition of electric Field Vector In Newtons Per Coulomb.
\end{proof}
\begin{definition}[electric Field Magnitude In Newtons Per Coulomb]
  \label{def:physics:phyx-mini-0833:phyxminiproblems-problemphyxmini0833-electricfieldmagnitudeinnewtonspercoulomb}
  \lean{PhyXMiniProblems.ProblemPhyXMini0833.electricFieldMagnitudeInNewtonsPerCoulomb}
  \uses{def:physics:phyx-mini-0833:phyxminiproblems-problemphyxmini0833-planarelectricfieldquantity, def:physics:phyx-mini-0833:phyxminiproblems-problemphyxmini0833-electricfieldvectorinnewtonspercoulomb}
  Euclidean magnitude of a planar electric field, in newtons per coulomb
\end{definition}
\begin{proof}
  This is the declared model definition of electric Field Magnitude In Newtons Per Coulomb.
\end{proof}
\begin{definition}[Source Charge]
  \label{def:physics:phyx-mini-0833:phyxminiproblems-problemphyxmini0833-sourcecharge}
  \lean{PhyXMiniProblems.ProblemPhyXMini0833.SourceCharge}
  The three source charges, named by their corners in the primary image
\end{definition}
\begin{proof}
  This is the declared model definition of Source Charge.
\end{proof}
\begin{definition}[Charge Sign]
  \label{def:physics:phyx-mini-0833:phyxminiproblems-problemphyxmini0833-chargesign}
  \lean{PhyXMiniProblems.ProblemPhyXMini0833.ChargeSign}
  The sign symbols shown inside the three coloured charge circles
\end{definition}
\begin{proof}
  This is the declared model definition of Charge Sign.
\end{proof}
\begin{definition}[Rectangle Edge]
  \label{def:physics:phyx-mini-0833:phyxminiproblems-problemphyxmini0833-rectangleedge}
  \lean{PhyXMiniProblems.ProblemPhyXMini0833.RectangleEdge}
  The four dashed sides making the rectangular layout visible
\end{definition}
\begin{proof}
  This is the declared model definition of Rectangle Edge.
\end{proof}
\begin{definition}[Separation Label]
  \label{def:physics:phyx-mini-0833:phyxminiproblems-problemphyxmini0833-separationlabel}
  \lean{PhyXMiniProblems.ProblemPhyXMini0833.SeparationLabel}
  The two printed separation labels and their geometric roles
\end{definition}
\begin{proof}
  This is the declared model definition of Separation Label.
\end{proof}
\begin{definition}[Rectangle Charge Figure]
  \label{def:physics:phyx-mini-0833:phyxminiproblems-problemphyxmini0833-rectanglechargefigure}
  \lean{PhyXMiniProblems.ProblemPhyXMini0833.RectangleChargeFigure}
  \uses{def:physics:phyx-mini-0833:phyxminiproblems-problemphyxmini0833-sourcecharge, def:physics:phyx-mini-0833:phyxminiproblems-problemphyxmini0833-chargesign, def:physics:phyx-mini-0833:phyxminiproblems-problemphyxmini0833-rectangleedge, def:physics:phyx-mini-0833:phyxminiproblems-problemphyxmini0833-separationlabel}
  Literal qualitative and textual content of image 833.png
\end{definition}
\begin{proof}
  This is the declared model definition of Rectangle Charge Figure.
\end{proof}
\begin{definition}[Rectangular Point Charge Setup]
  \label{def:physics:phyx-mini-0833:phyxminiproblems-problemphyxmini0833-rectangularpointchargesetup}
  \lean{PhyXMiniProblems.ProblemPhyXMini0833.RectangularPointChargeSetup}
  \uses{def:physics:phyx-mini-0833:phyxminiproblems-problemphyxmini0833-lengthquantity, def:physics:phyx-mini-0833:phyxminiproblems-problemphyxmini0833-signedchargequantity, def:physics:phyx-mini-0833:phyxminiproblems-problemphyxmini0833-planarpositionquantity, def:physics:phyx-mini-0833:phyxminiproblems-problemphyxmini0833-planarelectricfieldquantity, def:physics:phyx-mini-0833:phyxminiproblems-problemphyxmini0833-sourcecharge, def:physics:phyx-mini-0833:phyxminiproblems-problemphyxmini0833-rectanglechargefigure}
  The independent physical quantities in the electrostatic model. Each field contribution and the net field are observables, related only by the governing laws below; their values are not defined from an answer choice
\end{definition}
\begin{proof}
  This is the declared model definition of Rectangular Point Charge Setup.
\end{proof}
\begin{definition}[displacement From Source In Meters]
  \label{def:physics:phyx-mini-0833:phyxminiproblems-problemphyxmini0833-displacementfromsourceinmeters}
  \lean{PhyXMiniProblems.ProblemPhyXMini0833.displacementFromSourceInMeters}
  \uses{def:physics:phyx-mini-0833:phyxminiproblems-problemphyxmini0833-positionvectorinmeters, def:physics:phyx-mini-0833:phyxminiproblems-problemphyxmini0833-sourcecharge, def:physics:phyx-mini-0833:phyxminiproblems-problemphyxmini0833-rectangularpointchargesetup}
  Displacement in metres from a source charge to the marked field point
\end{definition}
\begin{proof}
  This is the declared model definition of displacement From Source In Meters.
\end{proof}
\begin{definition}[Matches Problem And Primary Figure]
  \label{def:physics:phyx-mini-0833:phyxminiproblems-problemphyxmini0833-matchesproblemandprimaryfigure}
  \lean{PhyXMiniProblems.ProblemPhyXMini0833.MatchesProblemAndPrimaryFigure}
  \uses{def:physics:phyx-mini-0833:phyxminiproblems-problemphyxmini0833-xaxis, def:physics:phyx-mini-0833:phyxminiproblems-problemphyxmini0833-yaxis, def:physics:phyx-mini-0833:phyxminiproblems-problemphyxmini0833-lengthinmeters, def:physics:phyx-mini-0833:phyxminiproblems-problemphyxmini0833-chargeincoulombs, def:physics:phyx-mini-0833:phyxminiproblems-problemphyxmini0833-positionvectorinmeters, def:physics:phyx-mini-0833:phyxminiproblems-problemphyxmini0833-rectangularpointchargesetup}
  Numerical and qualitative data transcribed from the problem and primary image. The coordinate origin is chosen at the marked lower-left point. No electric- field magnitude or answer value is included here
\end{definition}
\begin{proof}
  This is the declared model definition of Matches Problem And Primary Figure.
\end{proof}
\begin{definition}[Uses SIReference Data]
  \label{def:physics:phyx-mini-0833:phyxminiproblems-problemphyxmini0833-usessireferencedata}
  \lean{PhyXMiniProblems.ProblemPhyXMini0833.UsesSIReferenceData}
  \uses{def:physics:phyx-mini-0833:phyxminiproblems-problemphyxmini0833-rectangularpointchargesetup}
  Coherent-SI calibration of Physlib's Coulomb constant
\end{definition}
\begin{proof}
  This is the declared model definition of Uses SIReference Data.
\end{proof}
\begin{definition}[Has Physical Parameters]
  \label{def:physics:phyx-mini-0833:phyxminiproblems-problemphyxmini0833-hasphysicalparameters}
  \lean{PhyXMiniProblems.ProblemPhyXMini0833.HasPhysicalParameters}
  \uses{def:physics:phyx-mini-0833:phyxminiproblems-problemphyxmini0833-lengthinmeters, def:physics:phyx-mini-0833:phyxminiproblems-problemphyxmini0833-rectangularpointchargesetup, def:physics:phyx-mini-0833:phyxminiproblems-problemphyxmini0833-displacementfromsourceinmeters}
  Positivity and source/observation separation conditions
\end{definition}
\begin{proof}
  This is the declared model definition of Has Physical Parameters.
\end{proof}
\begin{definition}[Satisfies Point Charge Coulomb Law And Superposition]
  \label{def:physics:phyx-mini-0833:phyxminiproblems-problemphyxmini0833-satisfiespointchargecoulomblawandsuperposition}
  \lean{PhyXMiniProblems.ProblemPhyXMini0833.SatisfiesPointChargeCoulombLawAndSuperposition}
  \uses{def:physics:phyx-mini-0833:phyxminiproblems-problemphyxmini0833-chargeincoulombs, def:physics:phyx-mini-0833:phyxminiproblems-problemphyxmini0833-electricfieldvectorinnewtonspercoulomb, def:physics:phyx-mini-0833:phyxminiproblems-problemphyxmini0833-sourcecharge, def:physics:phyx-mini-0833:phyxminiproblems-problemphyxmini0833-rectangularpointchargesetup, def:physics:phyx-mini-0833:phyxminiproblems-problemphyxmini0833-displacementfromsourceinmeters}
  The point-charge law in vector form and linear superposition. For displacement r from source to observation point, a signed charge q contributes k q r / ‖r‖³. These are general physical relations and do not prescribe the requested magnitude or any displayed answer
\end{definition}
\begin{proof}
  This is the declared model definition of Satisfies Point Charge Coulomb Law And Superposition.
\end{proof}
\begin{definition}[Answer Choice]
  \label{def:physics:phyx-mini-0833:phyxminiproblems-problemphyxmini0833-answerchoice}
  \lean{PhyXMiniProblems.ProblemPhyXMini0833.AnswerChoice}
  Labels attached to the four displayed field-strength choices
\end{definition}
\begin{proof}
  This is the declared model definition of Answer Choice.
\end{proof}
\begin{definition}[displayed Field Strength In Newtons Per Coulomb]
  \label{def:physics:phyx-mini-0833:phyxminiproblems-problemphyxmini0833-displayedfieldstrengthinnewtonspercoulomb}
  \lean{PhyXMiniProblems.ProblemPhyXMini0833.displayedFieldStrengthInNewtonsPerCoulomb}
  \uses{def:physics:phyx-mini-0833:phyxminiproblems-problemphyxmini0833-answerchoice}
  Printed field strength beside each answer, in newtons per coulomb
\end{definition}
\begin{proof}
  This is the declared model definition of displayed Field Strength In Newtons Per Coulomb.
\end{proof}
\begin{definition}[Is Unique Closest Displayed Choice]
  \label{def:physics:phyx-mini-0833:phyxminiproblems-problemphyxmini0833-isuniqueclosestdisplayedchoice}
  \lean{PhyXMiniProblems.ProblemPhyXMini0833.IsUniqueClosestDisplayedChoice}
  \uses{def:physics:phyx-mini-0833:phyxminiproblems-problemphyxmini0833-electricfieldmagnitudeinnewtonspercoulomb, def:physics:phyx-mini-0833:phyxminiproblems-problemphyxmini0833-rectangularpointchargesetup, def:physics:phyx-mini-0833:phyxminiproblems-problemphyxmini0833-answerchoice, def:physics:phyx-mini-0833:phyxminiproblems-problemphyxmini0833-displayedfieldstrengthinnewtonspercoulomb}
  A displayed choice is uniquely closest to the calculated net-field magnitude. This comparison predicate is independent of which label ultimately satisfies it
\end{definition}
\begin{proof}
  This is the declared model definition of Is Unique Closest Displayed Choice.
\end{proof}
% --- Archon declaration topology end ---
% --- Archon physics formalization source end ---
